\documentclass[a4paper,12pt]{article}

\usepackage[T2A]{fontenc}
\usepackage[utf8]{inputenc}
\usepackage[russian]{babel}
\usepackage{amsmath,amssymb,amsthm,mathtools}
\usepackage{helvet}
\renewcommand{\familydefault}{\sfdefault}
\usepackage{icomma}
\usepackage[left=20mm,right=20mm,top=20mm,bottom=22mm]{geometry}
\usepackage{microtype}
\usepackage{tikz}
\usetikzlibrary{calc}
\usetikzlibrary{arrows.meta,positioning,shapes.geometric}
\usepackage{tkz-euclide}
\usepackage{pgfplots}
\pgfplotsset{compat=1.18}
\usepackage{tabularx,booktabs,longtable}
\usepackage{enumitem}
\usepackage{hyperref}
\usepackage{parskip}
\usepackage{fancyhdr}
\usepackage{setspace}
\usepackage{xcolor}

\definecolor{accent}{HTML}{008080}
\definecolor{darkaccent}{HTML}{004D4D}
\definecolor{textgray}{HTML}{333333}
\definecolor{softgray}{HTML}{E8EEEE}
\definecolor{linegray}{HTML}{879999}

\hypersetup{
    colorlinks=true,
    linkcolor=darkaccent,
    urlcolor=darkaccent,
    pdfauthor={Лёвин Артём Александрович},
    pdftitle={Параллелограмм в составных геометрических задачах}
}

\onehalfspacing
\setlength{\parindent}{0pt}
\setlength{\headheight}{15pt}
\setlist[itemize]{leftmargin=7mm,itemsep=2pt,topsep=3pt}
\setlist[enumerate]{leftmargin=9mm,itemsep=4pt,topsep=4pt}

\pagestyle{fancy}
\fancyhf{}
\fancyhead[L]{\small\color{textgray}Геометрия $\boldsymbol{\cdot}$ подготовка к ОГЭ}
\fancyhead[R]{\small\color{textgray}\thepage}
\renewcommand{\headrulewidth}{0.4pt}
\renewcommand{\headrule}{\hbox to\headwidth{\color{accent}\leaders\hrule height \headrulewidth\hfill}}

\tikzset{
    flowstart/.style={
        ellipse,
        draw=darkaccent,
        line width=0.9pt,
        minimum width=4.4cm,
        minimum height=1cm,
        align=center,
        fill=softgray
    },
    flowproc/.style={
        rectangle,
        rounded corners=3pt,
        draw=darkaccent,
        line width=0.9pt,
        text width=8.2cm,
        minimum height=1.15cm,
        align=center
    },
    flowcond/.style={
        diamond,
        aspect=2.3,
        draw=darkaccent,
        line width=0.9pt,
        text width=4.5cm,
        minimum height=1.25cm,
        align=center,
        inner sep=1.5pt,
        fill=softgray
    },
    flowarrow/.style={
        -{Stealth[length=3mm,width=2mm]},
        line width=0.9pt,
        draw=accent
    },
    flowlabel/.style={
        font=\small,
        text=darkaccent,
        fill=white,
        inner sep=1.5pt
    }
}

\tkzSetUpLine[line width=1pt,color=accent]
\tkzSetUpPoint[color=accent,fill=accent,size=3]
\tkzSetUpLabel[color=black,font=\small]

\begin{document}

\begin{titlepage}
\centering
\vspace*{2.2cm}

{\Large\color{darkaccent}\textbf{ГЕОМЕТРИЯ}}\\[1.2cm]

{\Huge\color{textgray}\textbf{Чек-лист}}\\[0.8cm]

{\LARGE\color{darkaccent}\textbf{Параллелограмм в составных задачах}}\\[0.35cm]

{\Large\color{textgray}Подобие треугольников, средние линии,\\
диагонали и теорема Вариньона}

\vfill

{\large София Кальней}\\[0.35cm]
{\large Подготовка к ОГЭ}\\[1.3cm]

{\large \today}\\[1.4cm]

{\large\color{darkaccent}
Лёвин Артём Александрович\\
эксклюзивно для Софии Кальней}

\vfill

\begin{tikzpicture}[remember picture,overlay]
\draw[accent,line width=1.2pt]
    ([xshift=2cm,yshift=1.55cm]current page.south west)
    .. controls
    ([xshift=6cm,yshift=2.25cm]current page.south west)
    and
    ([xshift=-6cm,yshift=0.85cm]current page.south east)
    ..
    ([xshift=-2cm,yshift=1.55cm]current page.south east);
\draw[linegray,line width=0.45pt]
    ([xshift=3cm,yshift=1.25cm]current page.south west)
    .. controls
    ([xshift=7cm,yshift=1.75cm]current page.south west)
    and
    ([xshift=-7cm,yshift=0.9cm]current page.south east)
    ..
    ([xshift=-3cm,yshift=1.25cm]current page.south east);
\end{tikzpicture}
\end{titlepage}

\section*{\color{darkaccent}1. Главная стратегия}

В составной геометрической задаче одна фигура часто скрывает несколько более простых конструкций. Основная цель --- выделить треугольники и установить связи между ними.

\begin{itemize}
    \item Выполните чертёж, отражающий форму фигур, взаимное расположение точек и порядок вершин.
    \item Выпишите данные: равенства, параллельность, середины сторон, углы и искомые величины.
    \item Найдите готовые треугольники или проведите диагонали.
    \item При наличии параллельных прямых проверьте подобие треугольников.
    \item Отметьте параллелограммы, прямоугольники и ромбы.
    \item Для площади выберите формулу, в которой известны почти все величины.
\end{itemize}

\subsection*{\color{darkaccent}Требования к чертежу}

\begin{itemize}
    \item Вершины многоугольника обозначаются последовательно по его границе.
    \item Ромб на рисунке должен иметь четыре визуально сопоставимые стороны.
    \item Точки, лежащие на сторонах, размещаются строго на соответствующих отрезках.
    \item Параллельность, прямые углы и равные отрезки отмечаются условными знаками.
    \item При неудобном чертеже допустимо изменить пропорции фигуры или порядок её расположения.
\end{itemize}

\section*{\color{darkaccent}2. Теоретическая опора}

\begin{table}[ht]
\centering
\caption{Свойства и формулы}
\renewcommand{\arraystretch}{1.35}
\begin{tabularx}{\linewidth}{>{\bfseries}p{0.28\linewidth}X}
\toprule
Объект & Ключевые факты \\
\midrule
Параллелограмм &
Противоположные стороны равны и параллельны; диагонали делятся точкой пересечения пополам. \\

Прямоугольник &
Все углы прямые; диагонали равны. \\

Ромб &
Все стороны равны; диагонали перпендикулярны, делят друг друга пополам и являются биссектрисами его углов. \\

Средняя линия треугольника &
Соединяет середины двух сторон, параллельна третьей стороне и равна её половине. \\

Подобные треугольники &
Соответственные углы равны, а соответственные стороны пропорциональны. Порядок вершин в записи подобия обязателен. \\

Прямоугольный треугольник &
Катет, лежащий против угла $30^\circ$, равен половине гипотенузы. \\

Выпуклый четырёхугольник &
Если диагонали равны $d_1$, $d_2$, а угол между ними равен $\varphi$, то
\[
S=\frac12 d_1d_2\sin\varphi.
\] \\
\bottomrule
\end{tabularx}
\end{table}

\subsection*{\color{darkaccent}Формулы площади параллелограмма}

Для сторон $a$, $b$, высоты $h_a$, угла $\alpha$ между сторонами и диагоналей $d_1$, $d_2$:

\[
S=ah_a,
\qquad
S=ab\sin\alpha,
\qquad
S=\frac12 d_1d_2\sin\varphi,
\]

где $\varphi$ --- угол между диагоналями.

\subsection*{\color{darkaccent}Тождество параллелограмма}

Если стороны параллелограмма равны $a$ и $b$, а диагонали --- $d_1$ и $d_2$, то

\[
\boxed{d_1^2+d_2^2=2(a^2+b^2)}.
\]

Формула позволяет найти неизвестную сторону или диагональ, когда известны остальные три величины.

\section*{\color{darkaccent}3. Ромб внутри прямоугольника}

В прямоугольнике $ABCD$ точки $K\in AB$ и $M\in CD$ выбраны так, что $AKCM$ --- ромб. Диагональ $AC$ образует со стороной $AB$ угол $30^\circ$. Пусть $AB=3$. Найдём сторону ромба.

\begin{center}
\begin{tikzpicture}[scale=0.9]
    \tkzDefPoint(0,0){A}
    \tkzDefPoint(6,0){B}
    \tkzDefPoint(6,3.464){C}
    \tkzDefPoint(0,3.464){D}
    \tkzDefPoint(4,0){K}
    \tkzDefPoint(2,3.464){M}

    \tkzDrawPolygon(A,B,C,D)
    \tkzDrawPolygon[thick,color=darkaccent](A,K,C,M)
    \tkzDrawSegment[dashed,color=linegray](A,C)

    \tkzDrawPoints(A,B,C,D,K,M)
    \tkzLabelPoints[below](A,K,B)
    \tkzLabelPoints[above](D,M,C)

    \tkzMarkRightAngle[size=0.28](A,B,C)
    \tkzMarkAngle[size=0.72,color=accent](K,A,C)
    \tkzLabelAngle[pos=1.05](K,A,C){$30^\circ$}

    \tkzLabelSegment[below](A,B){$3$}
\end{tikzpicture}
\end{center}

Диагональ ромба $AC$ делит углы при вершинах $A$ и $C$ пополам. Поэтому

\[
\angle KAC=\angle KCA=30^\circ.
\]

Поскольку $\angle BCD=90^\circ$ и точка $M$ лежит на $CD$, получаем

\[
\angle KCB=30^\circ.
\]

Треугольник $KBC$ прямоугольный. Обозначим $KB=x$. Катет $KB$ лежит против угла $30^\circ$, поэтому

\[
KC=2x.
\]

Все стороны ромба равны, следовательно, $AK=KC=2x$. Тогда

\[
AB=AK+KB=2x+x=3x.
\]

Из равенства $3x=3$ следует $x=1$. Сторона ромба равна

\[
\boxed{AK=KC=2}.
\]

\subsection*{\color{darkaccent}Что следует запомнить}

\begin{itemize}
    \item Информация о диагонали ромба сразу приводит к свойствам углов и перпендикулярности.
    \item После появления угла $30^\circ$ следует проверить наличие прямоугольного треугольника.
    \item Полезно вводить одну переменную для короткого катета.
\end{itemize}

\section*{\color{darkaccent}4. Параллелограмм внутри треугольника}

В треугольнике $ABC$ даны $AB=15$ и $AC=10$. Точки $D\in AB$, $E\in AC$, $P\in BC$ образуют параллелограмм $ADPE$. Известно, что $AD=6$. Найдём $AE$.

\begin{center}
\begin{tikzpicture}[scale=0.95]
    \tkzDefPoint(0,0){A}
    \tkzDefPoint(6,0){B}
    \tkzDefPoint(0,4){C}
    \tkzDefPoint(2.4,0){D}
    \tkzDefPoint(0,2.4){E}
    \tkzDefPoint(2.4,2.4){P}

    \tkzDrawPolygon(A,B,C)
    \tkzDrawPolygon[thick,color=darkaccent](A,D,P,E)
    \tkzDrawPoints(A,B,C,D,E,P)

    \tkzLabelPoints[below](A,D,B)
    \tkzLabelPoints[left](E,C)
    \tkzLabelPoints[right](P)

    \tkzLabelSegment[below](A,B){$15$}
    \tkzLabelSegment[left](A,C){$10$}
    \tkzLabelSegment[above](A,D){$6$}
\end{tikzpicture}
\end{center}

В параллелограмме $ADPE$ выполняется $DP\parallel AC$. Поэтому

\[
\triangle BDP\sim\triangle BAC.
\]

Имеем

\[
BD=AB-AD=15-6=9,
\qquad
\frac{BD}{BA}=\frac{9}{15}=\frac35.
\]

Из подобия:

\[
\frac{DP}{AC}=\frac35,
\qquad
DP=\frac35\cdot 10=6.
\]

Противоположные стороны параллелограмма равны, поэтому

\[
\boxed{AE=DP=6}.
\]

\subsection*{\color{darkaccent}Методический вывод}

При вписанном в треугольник параллелограмме:

\begin{enumerate}
    \item найдите пары параллельных прямых;
    \item выделите малый и большой подобные треугольники;
    \item строго сопоставьте соответственные стороны;
    \item используйте равенство противоположных сторон параллелограмма.
\end{enumerate}

\section*{\color{darkaccent}5. Теорема Вариньона}

Пусть $K$, $L$, $M$, $N$ --- середины сторон произвольного четырёхугольника $ABCD$.

\begin{center}
\begin{tikzpicture}[scale=1.05]
    \tkzDefPoint(0,0){A}
    \tkzDefPoint(5,0.7){B}
    \tkzDefPoint(4.2,4){C}
    \tkzDefPoint(-0.7,2.7){D}

    \tkzDefMidPoint(A,B)\tkzGetPoint{K}
    \tkzDefMidPoint(B,C)\tkzGetPoint{L}
    \tkzDefMidPoint(C,D)\tkzGetPoint{M}
    \tkzDefMidPoint(D,A)\tkzGetPoint{N}

    \tkzDrawPolygon(A,B,C,D)
    \tkzDrawSegments[dashed,color=linegray](A,C B,D)
    \tkzDrawPolygon[thick,color=darkaccent](K,L,M,N)

    \tkzDrawPoints(A,B,C,D,K,L,M,N)
    \tkzLabelPoints[below](A,K,B)
    \tkzLabelPoints[right](L,C)
    \tkzLabelPoints[above](M,D)
    \tkzLabelPoints[left](N)
\end{tikzpicture}
\end{center}

В треугольнике $ABC$ отрезок $KL$ является средней линией:

\[
KL\parallel AC,
\qquad
KL=\frac12 AC.
\]

Аналогично в треугольнике $ADC$:

\[
MN\parallel AC,
\qquad
MN=\frac12 AC.
\]

Следовательно, $KL=MN$ и $KL\parallel MN$. Аналогично

\[
LM=KN,
\qquad
LM\parallel KN.
\]

Таким образом,

\[
\boxed{KLMN\text{ --- параллелограмм}.}
\]

Это утверждение называется \textbf{теоремой Вариньона}.

\subsection*{\color{darkaccent}Связь с диагоналями исходного четырёхугольника}

Стороны параллелограмма Вариньона удовлетворяют равенствам

\[
KL=MN=\frac12 AC,
\qquad
LM=KN=\frac12 BD.
\]

Отсюда следуют полезные выводы:

\begin{itemize}
    \item если $AC\perp BD$, то $KLMN$ --- прямоугольник;
    \item если $AC=BD$, то $KLMN$ --- ромб;
    \item если $AC=BD$ и $AC\perp BD$, то $KLMN$ --- квадрат;
    \item периметр параллелограмма Вариньона равен
    \[
    P_{KLMN}=AC+BD.
    \]
\end{itemize}

\subsection*{\color{darkaccent}Отношение площадей}

Пусть $\varphi$ --- угол между диагоналями $AC$ и $BD$. Тогда

\[
S_{ABCD}=\frac12 AC\cdot BD\sin\varphi.
\]

Соседние стороны параллелограмма Вариньона равны половинам диагоналей исходного четырёхугольника, поэтому

\[
S_{KLMN}
=
\frac{AC}{2}\cdot\frac{BD}{2}\sin\varphi
=
\frac14 AC\cdot BD\sin\varphi.
\]

Следовательно,

\[
\boxed{S_{KLMN}=\frac12 S_{ABCD}}.
\]

\section*{\color{darkaccent}6. Обратная идея через параллелограмм Вариньона}

Отрезки, соединяющие середины противоположных сторон четырёхугольника, являются диагоналями параллелограмма Вариньона.

Если эти отрезки равны, то диагонали параллелограмма Вариньона равны. Значит, параллелограмм Вариньона является прямоугольником.

Его соседние стороны параллельны диагоналям исходного четырёхугольника. Поэтому диагонали исходного четырёхугольника перпендикулярны.

Если их длины равны $d_1$ и $d_2$, то

\[
\boxed{S=\frac12 d_1d_2}.
\]

\section*{\color{darkaccent}7. Типичные ошибки}

\begin{itemize}
    \item Вершины подобных треугольников записаны в несогласованном порядке.
    \item Равенство сторон параллелограмма используется до доказательства принадлежности фигуры к этому виду.
    \item Средней линией называют отрезок, хотя середины обеих сторон ещё не установлены.
    \item В формуле площади четырёхугольника пропускают множитель $\frac12$.
    \item В формуле через диагонали берут угол между сторонами.
    \item Свойство диагоналей прямоугольника или ромба переносят на произвольный параллелограмм.
    \item По чертежу делают вывод о равенстве или перпендикулярности без обоснования.
    \item В задаче с углом $30^\circ$ путают катет и гипотенузу.
\end{itemize}

\section*{\color{darkaccent}8. Итоговый чек-лист}

Перед завершением решения проверьте:

\begin{itemize}
    \item[$\square$] порядок вершин всех многоугольников;
    \item[$\square$] соответствие чертежа условию;
    \item[$\square$] наличие прямоугольных и подобных треугольников;
    \item[$\square$] использование всех условий параллельности;
    \item[$\square$] доказательство свойств вспомогательных фигур;
    \item[$\square$] правильный выбор формулы площади;
    \item[$\square$] множители $\frac12$ в формулах;
    \item[$\square$] единицы измерения и разумность ответа.
\end{itemize}

\clearpage
\thispagestyle{empty}



\clearpage

\section*{\color{darkaccent}Тренировочный блок}

\textbf{Путь А.} Задания расположены по нарастающей сложности. Решения оформляйте с обоснованием каждого свойства.

\subsection*{\color{darkaccent}Средний уровень}

\begin{enumerate}
    \item В прямоугольнике $ABCD$ точки $K\in AB$ и $M\in CD$ выбраны так, что $AKCM$ --- ромб. Диагональ $AC$ образует со стороной $AB$ угол $30^\circ$. Найдите сторону ромба, если $AB=12$.

    \item Стороны параллелограмма равны $7$ и $9$, одна из диагоналей равна $8$. Найдите вторую диагональ.

    \item Стороны параллелограмма равны $8$ и $5$, угол между ними равен $30^\circ$. Найдите площадь параллелограмма.
\end{enumerate}

\subsection*{\color{darkaccent}Уровень выше среднего}

\begin{enumerate}[resume]
    \item В треугольнике $ABC$ даны $AB=15$ и $AC=10$. Точки $D\in AB$, $E\in AC$, $P\in BC$ образуют параллелограмм $ADPE$. Найдите $AE$, если $AD=6$.

    \item Диагонали выпуклого четырёхугольника равны $10$ и $14$. Найдите периметр параллелограмма, вершинами которого служат середины сторон исходного четырёхугольника.

    \item Диагонали выпуклого четырёхугольника равны $16$ и $10$, угол между ними равен $30^\circ$. Найдите площадь параллелограмма Вариньона.

    \item Отрезки, соединяющие середины противоположных сторон выпуклого четырёхугольника, равны. Диагонали исходного четырёхугольника равны $18$ и $14$. Найдите его площадь.

    \item Диагонали ромба равны $12$ и $16$. Найдите сторону и площадь ромба.

    \item Стороны параллелограмма равны $5$ и $7$, одна из диагоналей равна $8$. Найдите вторую диагональ.
\end{enumerate}

\subsection*{\color{darkaccent}Сложный уровень}

\begin{enumerate}[resume]
    \item В выпуклом четырёхугольнике $ABCD$ диагонали равны $20$ и $14$. Точки $K$, $L$, $M$, $N$ являются серединами сторон $AB$, $BC$, $CD$, $DA$. Диагонали параллелограмма $KLMN$ равны. Найдите:
    \begin{enumerate}[label=\arabic*)]
        \item площадь четырёхугольника $ABCD$;
        \item периметр параллелограмма $KLMN$;
        \item площадь параллелограмма $KLMN$.
    \end{enumerate}
\end{enumerate}

\clearpage

\section*{\color{darkaccent}Ответы к тренировочному блоку}

\begin{table}[ht]
\centering
\caption{Краткие ответы}
\renewcommand{\arraystretch}{1.45}
\begin{tabularx}{\linewidth}{>{\centering\arraybackslash}p{1.2cm}X}
\toprule
№ & Ответ \\
\midrule
1 & $8$. \\

2 & $14$. \\

3 & $20$. \\

4 & $6$. \\

5 & $24$. \\

6 & $20$. \\

7 & $126$. \\

8 & Сторона $10$, площадь $96$. \\

9 & $2\sqrt{21}$. \\

10 & 1) $140$; \quad 2) $34$; \quad 3) $70$. \\
\bottomrule
\end{tabularx}
\end{table}

\vfill

\begin{center}
\color{darkaccent}
\textbf{Ключевая мысль занятия}

Сложный четырёхугольник становится понятнее после выделения треугольников,
проведения диагоналей и установления параллельности.
\end{center}

\end{document}